\section{Introduction}
\label{sec:intro}

Whether a video was really captured by a trustworthy device has become a practical question for newsrooms, courts and platforms. The industry answer is the C2PA standard~\cite{c2pa_spec}: the capture device signs a manifest that binds a hash of the asset to assertions about its origin, and the manifest travels with the file as metadata. This design has three properties that are undesirable in some settings. First, the manifest is \emph{separable}: social networks and many transcoding pipelines strip metadata, which is why C2PA added watermark-based ``soft bindings'' to recover a lost manifest. Second, it is \emph{identifying}: the signature is verified against a device or vendor certificate, so the verifier learns who signed. Third, it is \emph{visible}: anyone can see that the file carries provenance, and an adversary who wants to suppress it knows exactly which files to target.

Zero-knowledge proofs address the second point. A camera can prove that it belongs to a set of authorized devices without revealing which one, the pattern behind anonymous credentials and Semaphore-style group membership~\cite{zerocash2014,semaphore}. Recent work applies zk-SNARKs to media provenance, mostly to prove that a published image or video is a permissible transformation of a signed original~\cite{photoproof2016,zkimg2022,veritas2024,vimz2024,zkcinema2026,aipa2026}. Steganography addresses the first and third points: a message hidden in the coded bitstream survives metadata stripping and container changes, and is invisible to a party that does not hold the key~\cite{simmons1984,cachin1998,fridrich2009book}. Compressed-domain data hiding in H.264 has a long history, including embedding in the sign bits of CAVLC trailing ones, which leaves the bit rate unchanged~\cite{kapotas2007,mobasseri2007,fragile_h264_mtap}.

This paper combines the two. The question we answer is: \emph{can a short zero-knowledge proof of anonymous camera membership be hidden inside the H.264 bitstream it vouches for, bound to that exact bitstream, at a quality cost that can be predicted and controlled?} Three technical obstacles stand in the way. (i)~A proof embedded into a video changes the video, so a proof that commits to the video's hash cannot be computed before embedding; the binding has to be defined over a canonical form that embedding leaves unchanged, without opening a gap an attacker can exploit. (ii)~The proof must be small enough for a covert channel whose capacity is a few dozen bits per frame, and its verifier must not need the camera's secret---otherwise any verifier could forge proofs, and the proof would add nothing over a shared-key MAC. (iii)~Sign flips in intra-coded frames propagate through intra prediction, so the quality cost is not simply proportional to the number of embedded bits; it must be modeled before an embedding budget can be chosen.

\textbf{Contributions.}
\begin{itemize}
\item A \emph{video-bound anonymous camera proof} (Section~\ref{sec:design}): a Circom circuit of 8{,}737 constraints with three public inputs (registry root and a 256-bit binding split in two halves) proving that $\Pos(s)$ is a leaf of a depth-16 Poseidon Merkle registry, with Groth16 proofs of 129 bytes.
\item A \emph{carrier-masked video digest}: SHA-256 over the RBSP of every NAL unit with exactly the keyed carrier bits cleared. We prove that it is invariant under embedding and that every other edit is detected, either through the digest or through the extracted proof input (Section~\ref{sec:security}).
\item A \emph{keyed, segment-wise covert channel} on CAVLC trailing-one signs with HKDF-derived schedule and whitening keys, a three-byte frame, and one segment per IDR, which works identically on files and on live streams.
\item A \emph{closed-form distortion model} of trailing-one sign embedding, derived from the H.264 dequantization and inverse transform, together with its inverse: the number of flips, and hence the per-IDR bit cap, that keeps a frame above a target PSNR (Section~\ref{sec:distortion}).
\item An implementation and an evaluation (Sections~\ref{sec:impl}--\ref{sec:eval}): 27 encodes at three resolutions, a model validation on sequence-level and single-flip experiments under two encoder configurations, Groth16/PLONK costs, and an 11-case attack matrix.
\end{itemize}

\textbf{Scope.} The scheme is \emph{fragile} by design: it protects bitstreams that are transferred unchanged (evidence chains, archives, files sent as files). A re-encode destroys both the payload and the digest; robustness to transformations is the subject of the transformation-proof literature~\cite{photoproof2016,veritas2024,zkcinema2026} and is outside this paper. The proof states that a registered camera vouched for a bitstream, not that the scene is real.
